% Example "nonlinchem": the problem, the solutions against the exact one, and
% the convergence and invariant tests. Every number comes from data/
% (parameters.tex, results.tex, exact.dat, solution_<method>.dat,
% history_<method>.dat, error_<method>.dat), written by
% examples/nonlinchem/main.f90, and data/dp54_mpmath.dat, written by
% dp54_mpmath.py: run `fpm run --example nonlinchem` from the project root and
% `python3 dp54_mpmath.py` here, then build this file in examples/nonlinchem/.
\documentclass[11pt,a4paper]{article}
\usepackage[margin=2.2cm]{geometry}
\usepackage{amsmath,booktabs}
\usepackage{pgfplots}
\usepgfplotslibrary{groupplots}
\pgfplotsset{compat=1.18}
\pagestyle{empty}

\begin{document}
\section*{Example: nonlinear chemistry}

\subsection*{1. Problem}
From the SciML problem library (\texttt{prob\_ode\_nonlinchem}): a species
$y_1$ decays into $y_2$, which reacts with itself into $y_3$,
\begin{equation*}
  y_1' = -y_1, \qquad
  y_2' = y_1 - y_2^2, \qquad
  y_3' = y_2^2, \qquad
  y(0) = (1, 0, 0).
\end{equation*}
The exact solution: $y_1 = e^{-t}$; the equation for $y_2$ is of Riccati
type, and $y_2 = u'/u$ turns it into $u'' = e^{-t}u$, which in
$s = 2e^{-t/2}$ is the modified Bessel equation of order zero. With
$y_2(0) = 0$,
\begin{equation*}
  y_2(t) = -\frac{s}{2}\,
  \frac{I_1(s) - c\,K_1(s)}{I_0(s) + c\,K_0(s)},
  \qquad c = \frac{I_1(2)}{K_1(2)},
  \qquad y_3 = 1 - y_1 - y_2 ,
\end{equation*}
since $y_1 + y_2 + y_3$ is constant (the three rates add up to zero). The
term $y_2^2$ makes the problem nonlinear, and $y_2$ decays only like $1/t$.
\begin{center}
\small
% Written by examples/nonlinchem/main.f90: do not edit.
\begin{tabular}{ll}
\toprule
initial condition & $y(0) = (1, 0, 0)$ \\
interval & $0 \le t \le 20$ \\
steps, convergence & $\Delta t = 0.5, 0.25, \dots, 0.5/128$ (8 levels) \\
error & $\max_i |y_i - y_i^{\mathrm{exact}}|$ at $t = 20$ \\
observed order & slope of the two finest levels with errors above $10^{-12}$ \\
invariant & $\max_n |y_1 + y_2 + y_3 - 1|$ over every run of the study \\
figures & solutions at $\Delta t = 0.5$; error histories at $\Delta t = 0.25$ \\
methods & \texttt{rk3ws}, \texttt{rk4}, \texttt{dp54} \\
interface & plain arrays, \texttt{rk\_array\_scheme\_t} \\
exact solution & Bessel series (30 terms), within $1.55\times10^{-15}$ (relative) of mpmath at 30 digits \\
\bottomrule
\end{tabular}

\end{center}

\subsection*{2. Solutions}
\begin{center}
\begin{tikzpicture}
  \begin{groupplot}[
      group style={group size=2 by 1, horizontal sep=1.9cm}, % two panels side by side
      width=6.6cm, height=6.2cm,                           % size of each panel
      grid=major, grid style={black!12},                   % light grid
      xlabel={$t$}, xmin=0, xmax=20,                       % both panels against time
    ]
    % ---- (a) the three species ----
    \nextgroupplot[
      title={(a) $y_1, y_2, y_3$ at $\Delta t = 0.5$},     % panel title
      ymin=0, ymax=1.05,                                   % concentrations
      legend to name=chemlegend, legend columns=4,         % legend below both panels
      legend style={font=\small, /tikz/every even column/.append style={column sep=0.3cm}},
    ]
    \addplot[black, thick] table[x=t, y=y1] {data/exact.dat};                 % exact y1
    \addlegendentry{exact}
    \addplot[black, thick, forget plot] table[x=t, y=y2] {data/exact.dat};    % exact y2
    \addplot[black, thick, forget plot] table[x=t, y=y3] {data/exact.dat};    % exact y3
    \addplot[blue!80!black, only marks, mark=diamond*, mark size=1.5pt] table[x=t, y=y1] {data/solution_rk4.dat}; % RK4 y1
    \addlegendentry{\texttt{rk4}}
    \addplot[blue!80!black, only marks, mark=diamond*, mark size=1.5pt, forget plot] table[x=t, y=y2] {data/solution_rk4.dat}; % RK4 y2
    \addplot[blue!80!black, only marks, mark=diamond*, mark size=1.5pt, forget plot] table[x=t, y=y3] {data/solution_rk4.dat}; % RK4 y3
    \node[font=\small, anchor=west] at (axis cs:0.6,0.9)  {$y_1$}; % label of the decaying species
    \node[font=\small, anchor=west] at (axis cs:3,0.5)    {$y_2$}; % label of the intermediate
    \node[font=\small, anchor=west] at (axis cs:12,0.82)  {$y_3$}; % label of the product
    % ---- (b) error history ----
    \nextgroupplot[
      title={(b) error at $\Delta t = 0.25$},              % panel title
      ymode=log, ylabel={$\max_i |y_i - y_i^{\mathrm{exact}}|$}, % max-norm error, log scale
      ymin=1e-13, ymax=1,                                  % from round-off to order one
      legend pos=south west, legend columns=2,             % own legend, in the empty corner
      legend style={font=\scriptsize},                    % small
    ]
    \addplot[green!50!black, thick, dashed] table[x=t, y=error] {data/history_rk3ws.dat}; % order 2 here
    \addlegendentry{\texttt{rk3ws}}
    \addplot[blue!80!black, thick]   table[x=t, y=error] {data/history_rk4.dat};    % order 4
    \addlegendentry{\texttt{rk4}}
    \addplot[red!75!black, thick]    table[x=t, y=error] {data/history_dp54.dat};   % order 5
    \addlegendentry{\texttt{dp54}}
  \end{groupplot}
  \node[below=0.2cm] at (current bounding box.south)        % legend centred under both panels
    {\pgfplotslegendfromname{chemlegend}};
\end{tikzpicture}
\end{center}
(a) At $\Delta t = 0.5$ \texttt{rk4} is on the exact curves. (b) The error
grows while $y_1$ feeds $y_2$, then decays with the solution;
\texttt{rk3ws} (dashed) is far above \texttt{rk4}: on this nonlinear problem
it is of order 2, not 3.

\newpage
\subsection*{3. Tests: convergence order and the linear invariant}

\paragraph{Prediction.} (i) The problem is nonlinear, so each method
converges at its general order $p$, not at its order on linear problems:
\texttt{rk3ws} ($c = (0, \tfrac13, \tfrac12)$, $b = (0, 0, 1)$) satisfies
$b^{\mathsf T}c = \tfrac12$ and $b^{\mathsf T}Ac = \tfrac16$ but
$b^{\mathsf T}c^2 = \tfrac14 \ne \tfrac13$, a condition that only problems
with $f'' \ne 0$ (here $y_2^2$) see. (ii) Every Runge--Kutta method keeps
linear invariants: each stage adds $\Delta t\,a_{ij}k_j$, and every $k_j$ has
components summing to zero, so $y_1 + y_2 + y_3$ stays 1 to round-off at any
step size.

\paragraph{Result.}
\begin{center}
\small
% Written by examples/nonlinchem/main.f90: do not edit.
\begin{tabular}{llccccc}
\toprule
method & name & order & linear order & observed & $\max|y_1 + y_2 + y_3 - 1|$ \\
\midrule
RK3 of Wicker and Skamarock & \texttt{rk3ws} & 2 & 3 &  1.99 & $4.22\times10^{-15}$ \\
Classic RK4 & \texttt{rk4} & 4 & 4 &  4.03 & $8.44\times10^{-15}$ \\
Dormand-Prince 5(4) & \texttt{dp54} & 5 & 5 &  6.50 & $5.77\times10^{-15}$ \\
\midrule
\texttt{rk3ws}, pre- and post-processed & & & &  3.00 & $2.44\times10^{-15}$ \\
\bottomrule
\end{tabular}

\end{center}

\begin{center}
\begin{tikzpicture}
  \begin{loglogaxis}[
      width=10.5cm, height=7.5cm,                          % figure size
      xlabel={time step $\Delta t$},                       % x: step
      ylabel={error at $t = 20$},                          % y: max-norm global error
      ymin=1e-21, ymax=1e-1,                               % includes the 40-digit dp54 run
      grid=major, grid style={black!12},                   % light grid
      legend pos=outer north east, legend cell align=left, % legend right of the axes
      legend style={font=\small},                          % compact legend
      every axis plot/.append style={thick, mark size=1.5pt}, % same weight for all curves
    ]
    \addplot[green!50!black, dashed, mark=triangle] table[x=dt, y=error] {data/error_rk3ws.dat}; % slope 2 here
    \addlegendentry{\texttt{rk3ws}}
    \addplot[blue!80!black, mark=diamond*]   table[x=dt, y=error] {data/error_rk4.dat};      % slope 4
    \addlegendentry{\texttt{rk4}}
    \addplot[red!75!black, mark=pentagon*]   table[x=dt, y=error] {data/error_dp54.dat};     % slope 6.5 before round-off
    \addlegendentry{\texttt{dp54}}
    \addplot[red!75!black, dotted, mark=pentagon] table[x=dt, y=error] {data/dp54_mpmath.dat}; % 40 digits
    \addlegendentry{\texttt{dp54}, 40 digits}
    \addplot[green!50!black, densely dotted, mark=o] table[x=dt, y=error] {data/error_rk3ws_processed.dat}; % slope 3
    \addlegendentry{\texttt{rk3ws}, processed}
  \end{loglogaxis}
\end{tikzpicture}
\end{center}
Every method converges at its general order, \texttt{rk3ws} at 2, and the
invariant holds to round-off for all of them. \texttt{rk3ws} has
\emph{effective order} 3: its only third-order defect, the local error
$-\tfrac{h^3}{24}f''(f,f)$, equals $h^3\bigl(\tfrac{d}{dt}X - f'X\bigr)$ with
$X = -\tfrac1{24}f'f$, because $\tfrac{d}{dt}(f'f) = f''(f,f) + f'(f'f)$ along
any solution; so its global error is $h^2\bigl(X(u(t)) - \Phi(t,0)X(u_0)\bigr)
+ O(h^3)$. Starting from $u_0 + h^2 X(u_0)$ and correcting the result by
$-h^2X(u)$ removes that term: the processed \texttt{rk3ws} (last row, dotted
circles) converges at order 3. \texttt{dp54} shows a slope of
about $6.5$ in double precision: its fifth- and sixth-order error terms have
opposite signs here, so the error first falls faster than $\Delta t^5$. In
40-digit arithmetic (dotted, \texttt{dp54\_mpmath.py}) the error crosses zero
near $\Delta t \approx 0.03$ and then converges at order 5; in double
precision it reaches round-off before that regime.
\end{document}
